\section{Discussion}
\label{sec:discussion}

\subsection{Which Architecture Should a Practitioner Pick?}
The results support a simple decision rule for small models on latency-sensitive paths.
\begin{enumerate}
\item \emph{If the target has no known structure, keep it shallow and minimize operators.} On the linear and
curved tasks every topology reached the noise floor, so the cheapest one, a single wide hidden layer with three
operations, becomes the best choice. It trains fastest per step and has the lowest batch-1 latency.
\item \emph{If the target is compositional and its stages are known, a staged modular hybrid is worth its
overhead.} \HB cut chained-task excess error by \chainExcessRatio$\times$ at \stepRatioHB$\times$ the training cost per
step and \latOneRatioHB$\times$ the batch-1 latency. For offline or batched inference, that trade is
attractive. For per-sample loops, this is not the case.
\item \emph{Measure latency at the deployment batch size.} The ranking at batch 1 (operator-bound) and at batch
4096 (memory-bound) differ. A benchmark at one batch size can recommend the incorrect model for another.
\end{enumerate}

\ifcompact\else  % dropped in the 6-page build (summarized in Section V-B)
\subsection{Hybrids Converge Faster}
A less expected result is how much faster the hybrids converge. \HB reached the curved-task noise floor in
\curvedConvHB\ epochs, and \HA in 24, against \curvedConvSeq\ for single models. We hypothesize that the
additive multi-path output acts like a remaining groupage ~\cite{he2016resnet,lakshminarayanan2017simple}, where several
sub-networks fit different parts of the target at once, so no single narrow path is required to carry the whole
function. When training compute rather than inference latency is the bottleneck, this can outweigh the hybrids'
higher per-step cost.
\fi

\ifcompact\else  % dropped in the 6-page build (overlaps Section IV)
\subsection{Lessons for Small-Scale Benchmarking}
The pilot study shows how easily a benchmark can produce a confident, wrong
conclusion. Three practices caught it and are cheap enough for any student or practitioner to conduct. First,
include a \emph{null control}: a pair that should show no difference, such as our identical Sequential and
Modular networks. Our benchmarking served multiple roles, exposing the pilot's framework confound, catching CPU contention during
one training pass, and revealing module-call overhead that operation counts alone miss. Second, hold the framework and training loop fixed so that timing measures the model and not
the software stack~\cite{shi2016benchmarking}. Third, check accuracy against small baselines. A model's output that gives
the loss being equal to the target variance has not learned anything~\cite{lipton2019troubling}.
\fi

\subsection{Limitations and Threats to Validity}
\textbf{Scale.} The models are tiny ($\approx$385 parameters) and the tasks shown for testing are one-dimensional. This is
done deliberately, as it isolates the topology from overall capacity and makes the five-seed run cheap. It also means the
total latencies are dominated by per-operator overhead in a way that will shrink, though not vanish, for larger
models, where compute and traffic take over~\cite{sze2017efficient}. The ranking of architectures at
batch~4096, where arithmetic begins to matter, is a better guide to larger-scale behavior than batch~1.

\textbf{Hardware and runtime.} All timings come from one Apple M5 CPU running eager-mode PyTorch. Graph
compilers (TorchScript, \texttt{torch.compile}, ONNX Runtime) fuse small operators and would likely create narrowing in the
count penalty we observe. GPUs and microcontrollers would have very different overhead profiles. We
therefore present latency \emph{trends} across architectures, not total numbers to be quoted for other
platforms.\todo{optional: rerun latency\_sweep.py with torch.compile or on a Raspberry Pi / Colab GPU for a
second hardware point}

\textbf{Task design.} The chained task was written by the same author who designed \HB, so \HB's structural
match to it is an intentional best case for modular inductive bias, not an unbiased test. We do not claim
that staged modularity helps on unknown compositional problems, only that it \emph{can} when the
decomposition is known.

\textbf{Statistics.} Five seeds detect large effects but not small ones. We report Welch's $t$-tests
without multiple-comparison correction, so single and marginal results should be read as suggestive.
Hyperparameters (learning rate, epochs, batch size) were fixed for all models rather than tuned per
architecture, which may disadvantage architectures that give preference for other settings.
